%قبل از شروع مسئله ترجیحاً چیزی ننویسید. اما اگر خواستید بنویسید، در قالب لیستی (فقرات و شمارش) بنویسید.
%برای اضافه کردن نمره، بعدی از شروع‌مسئله، [x] را اضافه کنید و به جای x  نمره مورد نظر را قرار دهید. برای زیرمسئله‌ها هم به همین صورت است.
%\شروع‌مسئله[10]
\شروع‌مسئله

% صورت اصلی مسئله را در این‌جا بنویسید.

قصد داریم که روش آدرس‌دهی مستقیم را بر روی یک آرایه بزرگ انجام بدهیم. در هنگام شروع کار مقادیر نامعتبری بر روی حافظه گرفته شده برای آرایه قرار دارند و به دلیل حجم بالای آن، مقداردهی اولیه آنان برای ما به صرفه نیست. روشی برای پیاده‌سازی آدرس‌دهی مستقیم روی این آرایه بزرگ ارائه دهید. هر عضو درون آرایه، باید فضایی از $O(1)$ بگیرد و عملیات‌های درج، حذف و جست‌وجو هم در $O(1)$ قابل انجام باشند. همچنین توجه کنید که ساختن حالت اولیه داده‌ساختار هم باید از $O(1)$ باشد و حق مقداردهی اولیه خانه‌های آرایه را نداریم. صرفا با یک درخواست می‌توانیم حافظه مدنظر را برای اعمال بعدیمان رزرو کنیم و ممکن است پیش از شروع کار مقادیر نامعتبری روی آرایه باشند.

راهنمایی: از یک آرایه (یا پشته) کمکی استفاده کنید که بفهمید مقدار درون هر خانه آرایه بزرگ معتبر است یا نه.

\solution{

داده ساختار کمکی ما یک پشته به نام $S$ خواهد بود.

در ابتدا یک پشته خالی می‌سازیم و فضای مورد نیاز برای آرایه بزرگ را هم رزرو می‌کنیم اما کاری با مقداردهی اولیه آن نخواهیم داشت تا عملیات ما از $O(1)$ باشد. روش کار ما به این صورت خواهد بود که هر عضو در آرایه بزرگ در خود مقدار کلید و همچنین یک پوینتر به خانه از $S$ را نگاه می‌دارد. هر عضو در $S$ هم دقیقا به نقطه مشخص شده‌اش در آرایه بزرگ پوینتر دارد. بدین ترتیب از طریق تناظر درست پوینترها متوجه معتبر بودن هر خانه آرایه بزرگ خواهیم شد.

برای درج مقدار $x$ یک پوینتر به نام $y$ به $S$ پوش می‌کنیم که آدرس خانه $x$ ام آرایه بزرگ را دارد. در آرایه بزرگ و در خانه $x$ ام هم علاوه بر کلید مد نظر، پوینتر به خانه نظیر در $S$ یعنی $y$ قرار می‌دهیم.

برای جست‌وجوی وجود مقدار $x$ به خانه $x$ ام آرایه بزرگ می‌رویم. اگر این خانه پوینتری به نقطه‌ای از $S$ داشت که آن نقطه در $S$ هم پوینتری به خانه $x$ ام آرایه بزرگ داشت، یعنی این خانه معتبر است و واقعا مقداری در آن درج شده است. در غیر این صورت یعنی مقداری در آن درج نشده و داده نامعتبری از قبل در آن قرار گرفته است.

برای حذف، نگاه می‌کنیم که خانه $x$ ام آرایه بزرگ به چه مقداری در $S$ اشاره می‌کند. ابتدا آدرس این خانه را  در $temp$ ذخیره می‌کنیم. سپس بالایی ترین عنصر پشته را که مثلا به خانه $z$ ام آرایه بزرگ اشاره می‌کند از پشته پاپ کرده و مقدار درون آن را در خانه‌ای که آدرسش را در $temp$ ذخیره کرده بودیم قرار می‌دهیم تا پشته ما حفره نداشته باشد. حال به خانه $z$ ام آرایه بزرگ رفته و مقدار اشاره‌گر مربوط به پشته را در آن به $temp$ تغییر می‌دهیم. همه این عملیات ها در $O(1)$ قابل انجام خواهند بود.


توجه: پیاده سازی چنین سیستمی با زبان‌هایی که مستقیما از پوینتر و اشاره‌گر پشتیبانی نمی‌کنند کمی دشوار خواهد بود. برای همین اگر خواستید این سیستم را به شکل عملی پیاده‌سازی کنید توصیه می‌شود از  \lr{C} یا  \lr{C++} استفاده کنید. 


	
}






%%%%%%%%%%%%%%%%%%%%%%%%%%%%%%%%%%%%%%%%%%%%%%%%%%%%%
% برخی دستورات مفید برای تمرین‌ها. %
%%%%%%%%%%%%%%%%%%%%%%%%%%%%%%%%%%%%%%%%%%%%%%%%%%%%%

% فرمول‌های نمونه:
%$f(x) = \Theta(x)$
%$T(x, y) \leq \log(x) + 2^y + \binom{2n}{n}$
% {\tt code\_function}


% لیست موارد بدون شماره‌گذاری:
%\شروع{فقرات}
%    \فقره مورد اول 
%        \شروع{فقرات}
%            \فقره مورد اول داخلی 
%                \شروع{فقرات}
%                    \فقره اولین مورد خیلی داخلی 
%                    \فقره دومین مورد خیلی داخلی 
%                \پایان{فقرات}
%            \فقره مورد دوم داخلی 
%        \پایان{فقرات}
%    \فقره مورد دوم 
%\پایان{فقرات}

% لیست موارد شماره‌گذاری شده:
%\شروع{شمارش}
%    \فقره یک مورد 
%    \فقره یک مورد دیگر 
%    \فقره باز هم یک مورد دیگر
%\پایان{شمارش}

% برای نوشتن یک فرمول چند خطی که یک شماره بخوره. برای حالت‌های دیگر به سایت زیر مراجعه کنید:
% https://www.overleaf.com/learn/latex/Aligning_equations_with_amsmath
%\begin{align} 
%    \begin{split}
%        (x+y)^3 &= (x+y)^2(x+y) \\
%                &= (x^2+2xy+y^2)(x+y) \\
%                &= (x^3+2x^2y+xy^2) + (x^2y+2xy^2+y^3) \\
%                &= x^3+3x^2y+3xy^2+y^3
%    \end{split}                                 
%\end{align}

% نوشتن ماتریس و آرایه دوبعدی
%\begin{align}
%    A = 
%    \begin{bmatrix}
%        A_{11} & A_{21} \\
%        A_{21} & A_{22}
%    \end{bmatrix}
%\end{align}

% اضافه کردن تصویر با شرح. برای اضافه کردن تصویر بدون شرح، فقط دستور درج‌تصویر را استفاده کنید.
%\شروع{شکل}[h!]
%\تنظیم‌ازوسط
%  \درج‌تصویر[پهنا=0.5\پهنای‌متن]{img.jpg}
%\شرح{شرح تصویر}
%\برچسب{تصویر:۱}
%\پایان{شکل}
%\پایان{زیرمسائل}

%کد زیر را ببینید.
%%\lstinputlisting{source_filename.py}
%\begin{latin}
%\begin{lstlisting}
%# Import the modules
%import sys
%import random
%
%ans = True
%
%while ans:
%question = raw_input("Ask the magic 8 ball a question: (press enter to quit) ")
%
%answers = random.randint(1,8)
%
%if question == "":
%sys.exit()
%
%elif answers == 1:
%print "It is certain"
%
%elif answers == 2:
%print "Outlook good"
%
%elif answers == 3:
%print "You may rely on it"
%
%elif answers == 4:
%print "Ask again later"
%
%elif answers == 5:
%print "Concentrate and ask again"
%
%elif answers == 6:
%print "Reply hazy, try again"
%
%elif answers == 7:
%print "My reply is no"
%
%elif answers == 8:
%print "My sources say no"
%\end{lstlisting}
%\end{latin}
